\section{Combat Phase}
\subsection{Who Goes First?}

All movement and reactions have now been resolved. Reveal all remaining orders. The player who began the Movement Phase also begins the Combat Phase and performs the first action. This action may be either shooting or resolving an assault.

If a detachment can both shoot and participate in an assault, these two actions must be activated separately during the Combat Phase.

Once the action has been resolved, the opposing player performs an action. The players continue alternating until all shooting and assaults have been resolved. A player cannot pass an activation during the Combat Phase.

\subsection{Resolving Shooting}

To shoot, a detachment must have an order that allows it to do so and must not have been pinned during the current turn. Engaging in an assault, either directly or through a Consolidation move, also prevents a detachment from shooting.

A detachment may be activated to shoot as long as it is permitted to do so, even if it has no valid target.

Several actions allow a detachment to shoot, depending on its order:

\begin{description}
    \item[Shooting:] A normal shooting attack without any particular bonuses or penalties.

    \item[First Fire:] Detachments performing First Fire may reroll results of 1 on their To-Hit rolls, provided they have not moved.

    \item[Snap Fire:] Snap Fire suffers a $-1$ penalty to its To-Hit rolls.
\end{description}

\subsubsection{Selecting Targets}

A shooting attack must be directed at a valid target. Each attack die from each base may be assigned to a different target. You may therefore divide your fire between several targets or concentrate it against a single target.

You must declare how the attacks from every weapon and base are allocated before rolling any dice. If the detachment has several weapons, the firing player determines the order in which they are resolved.

Range, line of sight, and firing arc must be checked for every target.

\paragraph{Valid Targets}

A valid target may take one of the following forms:

\begin{itemize}
    \item An enemy base must be visible, within range, and within the weapon's firing arc.
    \item A structure must be visible, within range, and within the weapon's firing arc. The weapon must also be capable of damaging structures.
\end{itemize}

\subsubsection{Range}

Only weapons with a Range characteristic may shoot. You must therefore check whether the target is within range of the weapon. Range is measured from the centre of the firing base to the centre of the target base.

\subsubsection{Line of Sight}

You must also determine whether the firing base has line of sight to its target. A base has line of sight if a straight line can be drawn from the centre of the firing base to the centre of the target base without crossing an obstacle that blocks line of sight.

\paragraph{Obstacles}

If the line of sight between the firing base and its target passes through a terrain feature or another base, the line of sight may be blocked. Their Blocking Class is used to determine whether they interfere with the line of sight.

First, determine the total class of the firing base, the target base, and the obstacle. To determine a base's total class, add its Blocking Class to the Blocking Class of any terrain feature it occupies. For example, an infantry base with Class 1 positioned on a Class 3 hill has a total class of 4.

The following situations may then apply:

\begin{itemize}
    \item If the obstacle is of the same or a higher class than both the firing base and the target base, it affects the line of sight.

    \item If the obstacle is of a lower class than both the firing base and the target base, it does not affect the line of sight.

    \item If the obstacle is of a lower class than one of the two bases and of the same class as the other, it does not affect the line of sight.

    \item If the obstacle is of a lower class than one of the two bases but of a higher class than the other, it does not affect the line of sight if the higher-class base is the one closest to the edge of the obstacle, or if that base has a class greater than twice the class of the obstacle.
\end{itemize}

An \textit{Occulting} terrain feature or a base that affects a line of sight blocks it completely. \textit{Obstructing} terrain does not block line of sight but provides a cover save.

Bases belonging to the same detachment as the firing base or the target base are not considered obstacles and do not interfere with line of sight.

By definition, if one base has line of sight to another, the second base also has line of sight to the first.

\begin{figure}
    \centering
    \includegraphics[width=0.5\linewidth]{combat-a.png}
    \caption{\textbf{Example 1:} The line drawn between the centres of the two Infantry bases does not pass through the obstacle, so the two bases have line of sight to one another.
    \textbf{Example 2:} The line drawn between the centres of the two Infantry bases passes through the Class 3 obstacle, so the two bases do not have line of sight to one another.}
    \label{fig:CombatA}
\end{figure}

\begin{figure}
    \centering
    \includegraphics[width=0.5\linewidth]{combat-b.png}
    \caption{The line drawn between the centre of the Norn Queen (Class 4) and the centre of the Tactical Space Marine base (Class 1) passes through a Class 3 obstacle. It is therefore necessary to determine which base is closest to the obstacle:
    \textbf{Example 1:} The Norn Queen is closest to the obstacle, so the shooting attack is possible.
    \textbf{Example 2:} The Tactical Space Marine base is closest to the obstacle, so the shooting attack is not possible.}
    \label{fig:CombatB}
\end{figure}

\subsubsection{Firing Arc}

There is a difference between having line of sight to a base and being able to fire at it. The target must also be within the weapon's firing arc.

By default, a weapon's firing arc depends on the class of the firing base. Class 1 and 2 bases have a 360\textdegree{} firing arc, while Class 3 and higher bases have a 180\textdegree{} firing arc covering their front arc.

\subsubsection{Templates}

Weapons that use a template have ``Template'' listed in the \textit{Dice} column of their unit profile and possess the \textit{Template (X)} ability. The size of the template is specified in the weapon's abilities.

Firing arc and range are measured between the centre of the firing base and the centre of the template.

A base is hit by a template if the centre of its base is covered by the template, even if the base itself is not within line of sight.

Templates may be placed according to the following restrictions:

\paragraph{Circular Template}

When using a circular template, the centre of the template must be within range, within line of sight, and within the weapon's firing arc.

The template may cover bases that the firing base cannot see, provided it covers at least one valid target. If additional templates are used, each template must meet these requirements.

\paragraph{Rangeless Flame Template}

When using a rangeless flame template, the template may be placed in any position provided it lies entirely within the weapon's firing arc. The narrow tip of the template must be placed over the centre of the firing base.

The template must also cover at least one valid target.

\paragraph{Ranged Flame Template}

When using a ranged flame template, its narrow end must be within line of sight and the entire template must be within the weapon's firing arc, although the entire template does not need to be within line of sight.

The template's central axis must point towards the centre of the target base, with its narrow end facing the firing base. The template must also cover at least one valid target.

\subsubsection{To-Hit Roll}

Roll one die for each attack die listed in the weapon's profile. For a template weapon, instead roll one die for each base whose centre is covered by the template.

Each result equal to or greater than the weapon's To-Hit value inflicts one hit. This success threshold may be modified by certain abilities or actions. Regardless of any modifiers, a natural roll of 6 on a To-Hit roll is always successful.

\paragraph{Grouped Dice Rolls}

Where possible, attacks from several bases may be grouped against another group of bases. Instead of assigning each attack die to a specific target, a number of attack dice are assigned to the targeted group.

The following conditions must be met:

\begin{itemize}
    \item The targeted bases must belong to the same detachment and must be identical.
    \item Only attacks with identical To-Hit values and Armour Penetration values may be grouped together.
    \item Dice may only be grouped against bases from the same detachment that have identical cover saves and protection abilities.
\end{itemize}

The purpose is to create uniform groups in which every attack has the same chance of hitting and every target has the same resistance to damage.

For template attacks, hits may only be grouped if every affected base has been covered the same number of times. Otherwise, resolve the hits in separate groups according to how many times each base was covered.

Dice need not be grouped; attacks may instead be assigned to a specific target.

\subsubsection{Saving Throw}

When a base is hit, it may make a saving throw to avoid taking damage. A saving throw succeeds if the result of 1D6, modified by the weapon's Armour Penetration (AP), is equal to or greater than the base's Save value.

If the saving throw fails, the base loses one Wound. In most cases, this destroys the base, and it is removed from the battlefield.

\paragraph{Rear Fire}

Armour is generally thicker at the front than at the rear. When firing at a Class 3 or higher base from within its rear 180\textdegree{} arc, the attack receives an additional $-1$ AP modifier.

\subsubsection{Removing Casualties}

Casualties are removed in order, beginning with the base closest to the firing base and proceeding to the farthest. If two or more bases are equally distant, the player suffering the casualties chooses which base is destroyed.

If casualties are inflicted upon a detachment containing units with multiple Wounds, the number of casualties must be maximised. This applies even if the closest base is not the one that must suffer the damage.

\subsubsection{Shooting While Engaged in an Assault}

A base engaged in an assault may shoot if it is not pinned and has an order that permits it to fire.

It suffers a $-1$ penalty to the To-Hit rolls of all its weapons unless it is at least three classes higher than every opposing base. It may target either the enemy bases engaged with it or other valid targets.

If the assault was resolved earlier during the turn and the base survived, it may shoot without this penalty.

\subsubsection{Shooting into an Assault}

An assault may be targeted without restriction, either by standard shooting attacks or by template weapons. In both cases, attacks suffer a $-1$ penalty to their To-Hit rolls unless the target is at least three classes higher than every opposing base engaged with it.

For direct-fire attacks, each attack that misses must be rerolled against an opposing base in base-to-base contact with the original target. The firing player chooses which eligible opposing base is targeted by the rerolled attack.

For template attacks, To-Hit rolls are made normally against every base whose centre is covered by the template. If a base is covered but an opposing base engaged with it is not, that opposing base may be hit in the same manner as with a direct-fire attack.

Casualties must be distributed across as many separate duels as possible. They cannot be allocated in a way that leaves some bases fighting two against one while other bases are left without an opponent.

\subsection{Resolving an Assault}

An assault represents close combat between several detachments whose bases are in base-to-base contact. An assault does not represent only physical combat; it also includes the use of short-ranged weapons and grenades.

When you choose to resolve an assault, every base involved in the melee fights, even if the bases belong to several different detachments. A base engaged in an assault within 6~cm of another base engaged in an assault is considered part of the same assault, even if the two bases belong to entirely different detachments.

The player activating the assault activates all detachments involved in it, in the order of their choice. These detachments may use abilities that can be used during their assault activation.

Surviving detachments may subsequently be activated again for a shooting action if any of their bases are still able to fire.

\subsubsection{Assault Weapons}

Firing arcs do not apply during an assault. Assault weapons may always be used, regardless of the direction of their target.

\subsubsection{Simultaneous Assaults}

Assaults are considered simultaneous. The abilities of participating bases remain active until the assault has been completely resolved.

Therefore, even if a Medic or Leader is killed early in the assault, their abilities remain active until the end of the assault.

\begin{figure}
    \centering
    \includegraphics[width=0.5\linewidth]{combat-c.png}
    \caption{Bases A, B, and C are covered by the template and are therefore attacked with a $-1$ To-Hit penalty.
    Base D may be hit if the To-Hit roll against Base C fails. In that case, make another To-Hit roll against Base D, still with a $-1$ To-Hit penalty.}
    \label{fig:CombatC}
\end{figure}

\subsubsection{Identifying Combat Groups}

To resolve an assault, first determine which bases are fighting one another. Then create combat groups, each consisting of one base from one side fighting one or more opposing bases in contact with it.

In the simplest case, a combat group consists of two bases in contact with one another. However, a base may sometimes be engaged with several opponents, making the groups less clearly defined.

In this situation, the player who initiated the assault's resolution decides which bases fight one another, while distributing the bases as evenly as possible between the combat groups.

\begin{figure}
    \centering
    \includegraphics[width=0.5\linewidth]{combat-d.png}
    \caption{Base A is only in contact with Base B, so A and B must belong to the same combat group. The same applies to Bases D and E.
    Base C could belong to either group. The player who initiated the resolution of the assault chooses which group Base C joins.}
    \label{fig:CombatD}
\end{figure}

\subsubsection{Resolving Combat Groups}

The combat groups are now resolved. The player with the greatest number of bases in a group is referred to as the attacker and chooses the order in which the group's duels are resolved.

To resolve a duel, each player rolls 2D6 and adds the Assault Factor (AF) of their base. The player with the highest total wins. The losing base loses one Wound, with no saving throw permitted.

If the assault totals are tied, neither base loses any Wounds.

If opposing bases remain in the combat group after the duel, another duel is resolved. The next opponent fighting the surviving base receives an outnumbering bonus of 1D6, which is added to its AF. The remainder of the duel is resolved in the same manner as the initial duel.

Continue resolving duels until no further duel can be fought within the group. Each additional duel grants the outnumbering side a further cumulative bonus of 1D6 to its AF.

\begin{figure}
    \centering
    \includegraphics[width=0.5\linewidth]{combat-e.png}
    \caption{The Space Marine player has the greater number of bases and therefore determines the order of the duels.
    The first duel is fought between A and B, and A survives. The second duel is fought between A and C. C rolls 3D6 + AF, but A survives again. The final duel is then fought between A and D, with D rolling 4D6 + AF.}
    \label{fig:CombatE}
\end{figure}

\subsubsection{Resolving the Assault outcome}

Once all combat groups have been resolved, the assault must be concluded.

If one side has lost all its bases engaged in the assault, that side is declared the loser. Otherwise, count the number of Wounds lost by each side. The combat result is the difference between the number of Wounds lost by the two sides. The side that lost the greatest number of Wounds loses the assault.

If both sides lost the same number of Wounds, the player who activated the assault during the Combat Phase wins it.

Detachments that won the assault and were affected by it do not make a Morale Test. Detachments that lost the assault and were affected by it must make a Morale Test.

The Morale Test receives a modifier of 1 for every two points of combat result. If a detachment fails its Morale Test, it receives a Fall Back order and immediately makes a Fall Back move.

If a detachment already had a Fall Back order before the assault was resolved, it automatically makes a Fall Back move.

Detachments belonging to the side that lost the assault which were not affected by it, as defined by the Morale rules, or which passed their Morale Test, make a Regroup move.

This Regroup move is not performed if:

\begin{itemize}
    \item The combined remaining Wounds of the losing side's detachments are at least twice the combined remaining Wounds of the winning side's detachments.
    \item A base is at least three classes higher than all opposing bases.
    \item Less than half of the detachment's bases participated in the assault.
    \item The base is immobilised or has the \textit{Flyer} ability.
\end{itemize}

If the side that lost the assault meets any of these conditions, its bases do not move. In all cases, the bases are no longer considered engaged in an assault.

Detachments belonging to the winning side that do not have a Fall Back order may now make a Take Position move.

\paragraph{Regroup Move}

A Regroup move is up to 12~cm and requires the losing side to move as far away from its enemies as possible. If several directions are possible, the player who won the assault chooses the direction.

The bases may freely choose their facing at the end of the movement, except for non-agile Titans, which are forced backwards while retaining their facing.

If a base moves more than 12~cm away from every enemy base, it may stop moving and does not have to complete the full 12~cm movement.

\paragraph{Take Position Move}

A Take Position move allows a base to move up to 6~cm without engaging in an assault.

\paragraph{Regroup and Take Position Movement Rules}

Regroup and Take Position moves follow the standard movement rules but ignore the zones of control of detachments involved in the assault.

Every base in the detachment may perform this movement, whether or not it was engaged in the assault. Detachment coherency must still be maintained.

If maintaining coherency is impossible, some of the detachment's bases are destroyed. The owning player chooses which bases are destroyed.

\begin{figure}
    \centering
    \includegraphics[width=0.5\linewidth]{combat-f.png}
    \caption{The Space Marines suffer five casualties, while the Tyranids suffer two casualties and the Norn Queen loses one Wound. The Tyranids therefore win the assault by 2 points.}
    \label{fig:CombatF}
\end{figure}

\begin{figure}
    \centering
    \includegraphics[width=0.5\linewidth]{combat-g.png}
    \caption{The Space Marines then make a Regroup move, and the Tyranids make a Take Position move.}
    \label{fig:CombatG}
\end{figure}
